\section{Shortcuts for a locally connected hierarchy}\label{sec:hierarchy}

We require one matrix that controls average resistances on several related
boundaries.  A hierarchy specifies the boundaries on which this control is
required.  The matrix need only obey an upper bound on cuts; an upper bound
in the positive semidefinite order would be a stronger condition.

\begin{definition}[Hierarchy and its boundaries]\label{def:hierarchy}
Let $H=(V,E)$ be a finite loopless multigraph.  A \emph{hierarchy} on $V$
is a finite rooted tree of distinct nonempty vertex sets.  Its root is $V$,
its leaves are the singleton sets, and the children of every nonleaf form
a partition of that node into at least two sets.  For a nonroot node $A$,
write $A^*$ for its parent and define
\[
 P(A)=\delta_H(A),\qquad
 O(A)=E_H(A,A^*\setminus A).
\]
All edge sets here retain multiplicities.  Given an integer $h\ge2$, the
hierarchy is \emph{locally $h$-connected} if $H[A]$ is $h$-edge-connected
for every node $A$.  For singleton nodes this condition is vacuous.
\end{definition}

In this definition $|O(A)|\ge h$ for every nonroot $A$: it is the boundary
of the nonempty proper set $A$ in the $h$-edge-connected graph $H[A^*]$.
An edge belongs to $O(A)$ for exactly two nodes of the full hierarchy.
Indeed, if $C$ is the least node containing both endpoints, these two
nodes are the two children of $C$ containing its endpoints.  Consequently
an edge belongs to at most two parent boundaries in any subcollection.

Orient each edge arbitrarily and let $b_e$ be its incidence column.  As
usual, $L_H=\sum_{e\in E}b_eb_e^{\mathsf T}$.  A \emph{shortcut matrix}
for $H$ means a symmetric positive definite matrix $D$ on $\mathbb R^V$
such that
\begin{equation}\label{eq:hierarchy-shortcut}
 \mathbf 1_S^{\mathsf T}D\mathbf 1_S\le |\delta_H(S)|
 \qquad(\varnothing\ne S\subsetneq V).
\end{equation}
Its edge resistance is $R_D(e)=b_e^{\mathsf T}D^{-1}b_e$.
The restriction to proper nonempty cuts in
\eqref{eq:hierarchy-shortcut} is essential: imposing the same inequality
at $S=V$ would contradict positive definiteness.  We do not impose
$D\preceq L_H$.

The following estimate is the hierarchy input used in this paper.
It is a polynomial consequence of the Tree-CP theorem of Anari and
Oveis Gharan~\cite[Theorem~3.5, printed p.~22]{AOG}.  We give a complete
proof in Appendix~\ref{app:hierarchy-proof}, following their dual and
geometric method with parameters chosen for this weaker bound.

\begin{theorem}[Polynomial shortcut bound]\label{thm:hierarchy-power}
There is an absolute integer $q_H\ge2$ such that the following holds for
each integer $q\ge q_H$.  Let $H$ have a locally $q$-connected hierarchy,
and let $\mathcal M$ be any collection of its nonroot nodes satisfying
$|O(A)|\ge |P(A)|/q$.  Then $H$ has a shortcut matrix $D$ such that
\begin{equation}\label{eq:hierarchy-power}
 \frac{1}{|O(A)|}\sum_{e\in O(A)}R_D(e)\le q^{-31/40}
 \qquad(A\in\mathcal M).
\end{equation}
The cutoff is independent of the graph, the hierarchy depth and the
marked collection $\mathcal M$.
\end{theorem}
